\documentclass[11pt,a4paper]{article}

% ---------- Packages ----------
\usepackage[margin=1in]{geometry}
\usepackage{amsmath, amssymb, amsthm, mathtools}
\usepackage{microtype}
\usepackage{algorithm}
\usepackage{algpseudocode}
\usepackage[colorlinks=true,
            linkcolor=blue!50!black,
            urlcolor=blue!50!black,
            citecolor=blue!50!black]{hyperref}

% ---------- Theorem environments ----------
\theoremstyle{plain}
\newtheorem{theorem}{Theorem}[section]
\newtheorem{proposition}[theorem]{Proposition}
\newtheorem{lemma}[theorem]{Lemma}
\newtheorem{corollary}[theorem]{Corollary}

\theoremstyle{definition}
\newtheorem{definition}[theorem]{Definition}

\theoremstyle{remark}
\newtheorem{remark}[theorem]{Remark}
\newtheorem{example}[theorem]{Example}

% ---------- Math shortcuts ----------
\newcommand{\R}{\mathbb{R}}
\newcommand{\N}{\mathbb{N}}
\newcommand{\E}{\mathbb{E}}
\newcommand{\Var}{\operatorname{Var}}
\newcommand{\Cov}{\operatorname{Cov}}
\newcommand{\F}{\mathcal{F}}
\newcommand{\Prob}{\mathbb{P}}
\newcommand{\Q}{\mathbb{Q}}
\newcommand{\dd}{\mathop{}\!\mathrm{d}}
\newcommand{\eqdef}{\overset{\mathrm{def}}{=}}
\newcommand{\indic}{\mathbf{1}}

% ---------- Document metadata ----------
\title{Quasi-Monte Carlo for the European Call:\\
       Halton, Sobol, and RQMC}
\author{Quant Finance Portfolio --- Phase 2, Block 3.1}
\date{\today}

\begin{document}
\maketitle

\begin{abstract}
Building on the foundations established in Block~3.0, we apply
quasi-Monte Carlo to the European call pricer. To make QMC
non-trivially multi-dimensional we use the multi-step Euler
discretisation of Block~1.2.1 with $N = 20$ time steps, giving an
integration problem of dimension $d = 20$. We compare four
estimators: IID Monte Carlo (baseline from Block~1.2.1), QMC with
Halton, QMC with Sobol, and randomized QMC (RQMC) with Sobol plus
digital shifting. The empirical convergence rates measured on a
log-log slope plot are expected to confirm the theoretical
predictions: slope $\approx -0.5$ for IID, slope close to $-1$ for
Sobol RQMC at this dimension. The validation script also confirms
that RQMC delivers a usable half-width via $R$-replication averaging.
\end{abstract}

\tableofcontents

\section{The integration problem}
\label{sec:integration-problem}

The Euler-discretised European call price is
\begin{equation}
\E^\Q\bigl[e^{-rT} \max(S_T^{(N)} - K, 0)\bigr],
\end{equation}
where $S_T^{(N)} = S_0 \prod_{k=1}^N (1 + r h + \sigma \sqrt h \, Z_k)$
with $h = T/N$ and $Z_1, \ldots, Z_N$ i.i.d.\ standard normals. With
$N = 20$ this is a $20$-dimensional integral against a Gaussian
measure. Equivalently, it is a $20$-dimensional integral against
the uniform measure on $[0, 1)^{20}$ via the change of variables
$Z_k = \Phi^{-1}(u_k)$:
\begin{equation}
C^{\mathrm{Euler}} \;=\; \int_{[0, 1)^{20}}
e^{-rT} \max\bigl(S_T^{(N)}(\Phi^{-1}(u_1), \ldots, \Phi^{-1}(u_{20}))
                  - K,\; 0\bigr) \dd u.
\end{equation}
This is the integral that QMC will attack.

\paragraph{Why $N = 20$ rather than $N = 1$?} The exact GBM sampler
of Block~1.1 needs only one normal per path, so the European call
under exact GBM is a $1$-dimensional integral. QMC at $d = 1$ is
trivial (van der Corput) and not pedagogically rich. Choosing the
Euler discretisation of Block~1.2.1 with $N = 20$ gives a
genuine $20$-dimensional problem where QMC's advantage matters.

The difference between exact and Euler is the $O(h)$ weak bias
documented in Block~1.2.1; for $N = 20$ this bias is approximately
$0.05$ in absolute terms, large enough to be visible against MC
error at $n \sim 10^4$ but small enough not to obscure the
convergence-rate analysis. We accept this bias as a fixed offset
and study the rate at which the Monte Carlo / QMC error around
$C^{\mathrm{Euler}}$ vanishes as $n$ grows.

\section{The four estimators}
\label{sec:estimators}

We construct four estimators of $C^{\mathrm{Euler}}$, all sharing
the same $\Phi^{-1}$ inversion step but differing in how the
underlying $u_i \in [0, 1)^{20}$ points are produced.

\paragraph{IID-MC.} Standard Monte Carlo. The points $u_i$ are
pseudo-random uniform draws from a Mersenne-Twister generator. This
is the Block~1.2.1 baseline.

\paragraph{QMC-Halton.} The points $u_i$ are the first $n$ entries
of the $20$-dimensional Halton sequence with primes $p_1 = 2,
p_2 = 3, \ldots, p_{20} = 71$. The estimator returns a single number
without a half-width.

\paragraph{QMC-Sobol.} The points $u_i$ are the first $n$ entries
of the $20$-dimensional Sobol sequence with Joe-Kuo direction
numbers. As with Halton, the deterministic estimator has no
half-width.

\paragraph{RQMC-Sobol.} The points are the Sobol sequence
randomised by digital shifting with $R$ independent shifts (we use
$R = 20$). Each replication produces an estimate of dimension
$d = 20$ from $n$ Sobol points; the $R$ replication estimates are
i.i.d.\ and a half-width is constructed from their sample variance.

\subsection{The fixed and varying axes}

To compare convergence rates, we vary the per-replication path
count $n$ across a logarithmic grid. The rest of the configuration
is fixed:
\begin{itemize}
\item Euler steps: $N = 20$.
\item Total dimensionality: $d = 20$.
\item RQMC replications: $R = 20$.
\item Random seed for IID and for RQMC's shifts: fixed for
reproducibility.
\item Target: the same canonical parameters $S_0 = K = 100$, $r =
0.05$, $\sigma = 0.20$, $T = 1$.
\end{itemize}
The reference value against which we measure error is the
closed-form Black-Scholes price $C^{\mathrm{BS}} = 10.4506$. The
Euler scheme's $O(h)$ bias creates a systematic offset of
approximately $0.05$ from $C^{\mathrm{BS}}$; this offset is the
\emph{same for all four estimators} (they all share the same
discretised payoff structure), so when we plot ``error vs $n$'' on
a log-log scale, the bias contributes a horizontal asymptote that
all four lines share. The slopes at moderate $n$, before the
asymptote dominates, are what we are after.

\section{Convergence rates: predictions}
\label{sec:rates}

The theoretical asymptotic rates from Block~3.0:
\begin{itemize}
\item IID-MC: error $= O(n^{-1/2})$. Log-log slope $-0.5$.
\item QMC-Halton: error $= O((\log n)^{20} / n)$. The $(\log n)^{20}$
factor at $n = 10^4$ is $(9.2)^{20} \approx 10^{19}$, so the constant
in the big-O is huge and Halton may not yet be in its asymptotic
regime at the budgets we test. Empirically, we expect a slope between
$-0.7$ and $-0.9$ rather than the asymptotic $-1$.
\item QMC-Sobol (deterministic): same asymptotic rate as Halton,
but with substantially better constants thanks to the $\mathbb{F}_2$
construction. Empirically, we expect a slope close to $-1$ at
moderate $n$.
\item RQMC-Sobol: same expected slope as deterministic Sobol; the
randomisation provides the half-width without affecting the rate.
\end{itemize}

\subsection{The dimension-effect caveat}

A naive reading of $(\log n)^d / n$ at $d = 20$ would predict QMC
to require astronomical $n$ to outperform MC. The empirical
behaviour is much friendlier, for two reasons:

\paragraph{Effective dimension.} The European call payoff, expressed
in terms of $u \in [0, 1)^{20}$, depends primarily on a small
number of ``effective dimensions'': the few coordinates whose
variation drives most of the variance of $S_T^{(N)}$. For the Euler
discretisation, all $20$ steps contribute roughly equally to
$\log S_T^{(N)}$, so the effective dimension is closer to $20$ than
to $1$, but the integrand's smoothness ($\max(\cdot, 0)$ has a
single kink in $\R^{20}$) helps QMC's convergence in practice. The
``intrinsic'' difficulty is lower than the nominal $d = 20$ would
suggest.

\paragraph{Brownian bridge construction.} A standard QMC
optimisation in finance is to construct the path of $W_t$ in the
\emph{Brownian bridge} order rather than the natural left-to-right
order: assign the first dimension of $u$ to $W_T$ (the most
``important'' coordinate), the second to $W_{T/2}$, the third and
fourth to $W_{T/4}$ and $W_{3T/4}$, and so on. This concentrates the
``effective dimension'' in the first few QMC coordinates, where
QMC's discrepancy advantage is strongest. \emph{We do not apply
the Brownian bridge construction in this writeup}: the natural
left-to-right Euler discretisation already gives observable QMC
advantage and the bridge optimisation can be added in a future
block (it is most relevant for path-dependent options where the
effective dimension can be much larger).

\subsection{The convergence rate test}

The validation script computes the absolute error
\begin{equation}
\mathrm{err}(n) \;=\; \bigl|\hat C_n - C^{\mathrm{Euler}}\bigr|
\end{equation}
on a logarithmic grid $n \in \{10^3, 10^{3.5}, 10^4, 10^{4.5}, 10^5\}$
(five points), fits a linear regression of $\log \mathrm{err}$ on
$\log n$, and reports the slope. We use the bias-corrected target
$C^{\mathrm{Euler}}$ (estimated from a high-precision IID-MC run)
rather than $C^{\mathrm{BS}}$ to avoid contaminating the rate
measurement with the Euler bias.

\section{Implementation overview}
\label{sec:impl}

The Python implementation introduces a new module \texttt{quantlib.qmc}:

\begin{itemize}
\item \texttt{halton(n, d)}: returns the first $n$ entries of the
$d$-dimensional Halton sequence as an $(n, d)$ NumPy array.
Implemented from scratch using the radical inverse function.
\item \texttt{sobol(n, d, scramble=False, seed=None)}: returns the
first $n$ Sobol points. Wraps \texttt{scipy.stats.qmc.Sobol}, which
uses Joe-Kuo direction numbers internally.
\item \texttt{mc\_european\_call\_euler\_qmc}: deterministic QMC
pricer that returns a single estimate (no half-width). Accepts a
parameter \texttt{sequence} $\in$ \{\texttt{'halton'},
\texttt{'sobol'}\}.
\item \texttt{mc\_european\_call\_euler\_rqmc}: randomized QMC
pricer with $R$ digital shifts. Returns a full \texttt{MCResult}
with the half-width estimated from the $R$ replications.
\end{itemize}

The pricers reuse the Euler step from \texttt{quantlib.gbm}: only
the source of normals changes (from
\texttt{rng.standard\_normal(...)} to $\Phi^{-1}$ applied to a QMC
sequence). This is exactly the modular separation foreseen in
Block~0 when we adopted inversion-based normal sampling.

\section{Empirical validation strategy}
\label{sec:validation}

The validation script \texttt{python/validate\_mc\_european\_qmc.py}
runs four tests.

\paragraph{Test 1: convergence rate slopes.} Compute the absolute
error of each estimator (IID, Halton, Sobol-det, Sobol-RQMC) at
$n \in \{10^3, 10^{3.5}, 10^4, 10^{4.5}, 10^5\}$. Linear-regress
$\log \mathrm{err}$ on $\log n$. Predicted slopes:
\begin{itemize}
\item IID: $\approx -0.5$.
\item Halton: $\approx -0.7$ to $-0.9$.
\item Sobol-det: $\approx -1.0$ (or slightly less).
\item Sobol-RQMC: $\approx -1.0$ (same as deterministic Sobol; the
randomisation only adds a small constant overhead).
\end{itemize}

\paragraph{Test 2: BS-coherence of RQMC-Sobol.} The RQMC pricer at
$n = 10^4$ paths and $R = 20$ replications agrees with $C^{\mathrm{BS}}$
modulo the Euler bias, within a few half-widths.

\paragraph{Test 3: half-width below MC at equal budget.} RQMC-Sobol
at $n_{\text{paths}} \times R = 100\,000$ payoffs has half-width
significantly below the IID-MC half-width at the same total budget.

\paragraph{Test 4: input validation.} Mechanical: the QMC and RQMC
pricers reject bad inputs.

\section{Bibliographic notes}
\label{sec:biblio}

The application of QMC to financial option pricing is reviewed in
Glasserman~\cite[Chapter~5]{Glasserman}. The Brownian bridge
construction in QMC is due to Caflisch, Morokoff and Owen
\cite{CaflischMorokoffOwen}; we do not implement it here. The
analysis of effective dimension in finance is in Wang and Sloan
\cite{WangSloan}.

\begin{thebibliography}{99}

\bibitem{CaflischMorokoffOwen}
R.\ E.\ Caflisch, W.\ Morokoff, A.\ B.\ Owen,
\emph{Valuation of mortgage-backed securities using Brownian
bridges to reduce effective dimension},
Journal of Computational Finance, 1(1):27--46, 1997.

\bibitem{Glasserman}
P.\ Glasserman,
\emph{Monte Carlo Methods in Financial Engineering},
Springer, 2004.

\bibitem{WangSloan}
X.\ Wang and I.\ H.\ Sloan,
\emph{Why are high-dimensional finance problems often of low
effective dimension?},
SIAM Journal on Scientific Computing, 27(1):159--183, 2005.

\end{thebibliography}

\end{document}
